\documentclass[10pt,a4paper]{amsart}
\usepackage{amsmath,amssymb,amsthm,graphicx}
\usepackage{verbatim}
\newenvironment{reference}{\comment}{\endcomment}
%\newenvironment{reference}{}{}
\newenvironment{slogan}{\comment}{\endcomment}
\newenvironment{history}{\comment}{\endcomment}

% For commutative diagrams we use Xy-pic
\usepackage[all]{xy}

% We use 2cell for 2-commutative diagrams.
\xyoption{2cell}
% Theorem environments.
%
\theoremstyle{plain}
\newtheorem{theorem}[subsection]{Theorem}
\newtheorem{proposition}[subsection]{Proposition}
\newtheorem{lemma}[subsection]{Lemma}

\theoremstyle{definition}
\newtheorem{definition}[subsection]{Definition}
\newtheorem{example}[subsection]{Example}
\newtheorem{exercise}[subsection]{Exercise}
\newtheorem{situation}[subsection]{Situation}

\theoremstyle{remark}
\newtheorem{remark}[subsection]{Remark}
\newtheorem{remarks}[subsection]{Remarks}

\numberwithin{equation}{subsection}

% Macros
%
\def\lim{\mathop{\mathrm{lim}}\nolimits}
\def\colim{\mathop{\mathrm{colim}}\nolimits}
\def\Spec{\mathop{\mathrm{Spec}}}
\def\Hom{\mathop{\mathrm{Hom}}\nolimits}
\def\Ext{\mathop{\mathrm{Ext}}\nolimits}
\def\SheafHom{\mathop{\mathcal{H}\!\mathit{om}}\nolimits}
\def\SheafExt{\mathop{\mathcal{E}\!\mathit{xt}}\nolimits}
\def\Sch{\mathit{Sch}}
\def\Mor{\mathop{\mathrm{Mor}}\nolimits}
\def\Ob{\mathop{\mathrm{Ob}}\nolimits}
\def\Sh{\mathop{\mathit{Sh}}\nolimits}
\def\NL{\mathop{N\!L}\nolimits}
\def\CH{\mathop{\mathrm{CH}}\nolimits}
\def\proetale{{pro\text{-}\acute{e}tale}}
\def\etale{{\acute{e}tale}}
\def\QCoh{\mathit{QCoh}}
\def\Ker{\mathop{\mathrm{Ker}}}
\def\Im{\mathop{\mathrm{Im}}}
\def\Coker{\mathop{\mathrm{Coker}}}
\def\Coim{\mathop{\mathrm{Coim}}}

% Boxtimes
%
\DeclareMathSymbol{\boxtimes}{\mathbin}{AMSa}{"02}

%
% Macros for moduli stacks/spaces
%
\def\QCohstack{\mathcal{QC}\!\mathit{oh}}
\def\Cohstack{\mathcal{C}\!\mathit{oh}}
\def\Spacesstack{\mathcal{S}\!\mathit{paces}}
\def\Quotfunctor{\mathrm{Quot}}
\def\Hilbfunctor{\mathrm{Hilb}}
\def\Curvesstack{\mathcal{C}\!\mathit{urves}}
\def\Polarizedstack{\mathcal{P}\!\mathit{olarized}}
\def\Complexesstack{\mathcal{C}\!\mathit{omplexes}}
% \Pic is the operator that assigns to X its picard group, usage \Pic(X)
% \Picardstack_{X/B} denotes the Picard stack of X over B
% \Picardfunctor_{X/B} denotes the Picard functor of X over B
\def\Pic{\mathop{\mathrm{Pic}}\nolimits}
\def\Picardstack{\mathcal{P}\!\mathit{ic}}
\def\Picardfunctor{\mathrm{Pic}}
\def\Deformationcategory{\mathcal{D}\!\mathit{ef}}
\usepackage{cleveref}
\crefname{originalreference}{original reference}{original references}
\Crefname{originalreference}{Original reference}{Original references}
\makeatletter
\newlabel{situation-setup}{{02QL}{source}{}{sourcepointer-1}{}}
\newlabel{situation-setup@cref}{{[originalreference][0][]02QL}{[1][0][]source}{}{}{}}
\makeatother
\crefname{originallemma}{original lemma}{original lemmas}
\Crefname{originallemma}{Original lemma}{Original lemmas}
\makeatletter
\newlabel{lemma-compose-flat-pullback}{{02RD}{source}{}{sourcepointer-2}{}}
\newlabel{lemma-compose-flat-pullback@cref}{{[originallemma][0][]02RD}{[1][0][]source}{}{}{}}
\makeatother
\crefname{originallemma}{original lemma}{original lemmas}
\Crefname{originallemma}{Original lemma}{Original lemmas}
\makeatletter
\newlabel{divisors-lemma-section-smooth-regular-immersion}{{067R}{source}{}{sourcepointer-3}{}}
\newlabel{divisors-lemma-section-smooth-regular-immersion@cref}{{[originallemma][0][]067R}{[1][0][]source}{}{}{}}
\makeatother
\crefname{originallemma}{original lemma}{original lemmas}
\Crefname{originallemma}{Original lemma}{Original lemmas}
\makeatletter
\newlabel{divisors-lemma-flat-base-change-regular-immersion}{{067P}{source}{}{sourcepointer-4}{}}
\newlabel{divisors-lemma-flat-base-change-regular-immersion@cref}{{[originallemma][0][]067P}{[1][0][]source}{}{}{}}
\makeatother
\crefname{originallemma}{original lemma}{original lemmas}
\Crefname{originallemma}{Original lemma}{Original lemmas}
\makeatletter
\newlabel{lemma-composition-regular-immersion}{{0FF0}{source}{}{sourcepointer-5}{}}
\newlabel{lemma-composition-regular-immersion@cref}{{[originallemma][0][]0FF0}{[1][0][]source}{}{}{}}
\makeatother
\crefname{originalreference}{original reference}{original references}
\Crefname{originalreference}{Original reference}{Original references}
\makeatletter
\newlabel{remark-restriction-bivariant}{{0F9Z}{source}{}{sourcepointer-6}{}}
\newlabel{remark-restriction-bivariant@cref}{{[originalreference][0][]0F9Z}{[1][0][]source}{}{}{}}
\makeatother
\crefname{originallemma}{original lemma}{original lemmas}
\Crefname{originallemma}{Original lemma}{Original lemmas}
\makeatletter
\newlabel{lemma-construction-gysin}{{0FBK}{source}{}{sourcepointer-7}{}}
\newlabel{lemma-construction-gysin@cref}{{[originallemma][0][]0FBK}{[1][0][]source}{}{}{}}
\makeatother
\crefname{originallemma}{original lemma}{original lemmas}
\Crefname{originallemma}{Original lemma}{Original lemmas}
\makeatletter
\newlabel{lemma-section-smooth}{{0FF1}{source}{}{sourcepointer-8}{}}
\newlabel{lemma-section-smooth@cref}{{[originallemma][0][]0FF1}{[1][0][]source}{}{}{}}
\makeatother
\title{The Gysin class of a local complete intersection morphism is independent of the chosen factorization}
\author{The Stacks Project Authors}
\date{stacks-project:chow.tex:local-snapshot-20261012; GNU FDL 1.2 or later}
\begin{document}
\maketitle
\noindent\textit{Editorial note.} All delivered mathematical text is the author's own, copied from the pinned Stacks Project chapter file chow.tex; no statement or proof step was rewritten, reordered or removed.\par\medskip
\noindent\textit{Editorial note.} Preamble provenance: the pinned chapter file opens with an input of the shared preamble and carries no preamble text of its own, so the project's single Stacks-wide preamble staged at sources/stacks-chow/preamble.tex is excerpted verbatim as two spans, lines 14-18 (Xy-pic setup) and lines 154-217 (the native theorem declarations and mathematical macros). Nothing was added, removed or redefined; the wrapper class and the amsmath/amssymb/amsthm/graphicx base packages are supplied by the assembler.\par\medskip
\noindent\textit{Editorial note.} The counted statement and proof contain no cross reference, no citation and no included file, so the delivered body is self contained.\par\medskip
\noindent\textit{Editorial note.} Canonical location of the counted proposition: Stacks tag 0FF2 (sources/stacks-chow/pinned-original-source.tex line 13485).\par\medskip
\noindent\textit{Allowed standard prerequisites (editorial).}\begin{itemize}
\item Composition of Gysin classes and the Gysin class of a regular immersion. Same chapter, Lemmas labelled lemma-composition-regular-immersion and lemma-construction-gysin, cited at proof lines 13517-13523. Splits (i')\textasciicircum{}! = s\textasciicircum{}! o j\textasciicircum{}! and identifies j\textasciicircum{}! as a restriction of i\textasciicircum{}!.
\item Regular immersions obtained from a section of a smooth morphism and from flat base change. Divisors, Lemma labelled divisors-lemma-section-smooth-regular-immersion and Lemma labelled divisors-lemma-flat-base-change-regular-immersion, cited at proof lines 13513-13517. Makes s and j regular immersions so that their Gysin classes exist.
\item Composition of flat pullbacks and commutation of bivariant classes with flat pullback. Same chapter, Lemma labelled lemma-compose-flat-pullback, cited at proof lines 13505-13507, and the flat pullback commutation used at proof lines 13523-13525. Splits (g')\textasciicircum{}* and moves the flat pullback past j\textasciicircum{}!.
\item The Gysin class of a section of a smooth morphism is the identity. Same chapter, Lemma labelled lemma-section-smooth, cited at proof lines 13525-13527. Finishes the reduction with s\textasciicircum{}! o p-bar\textasciicircum{}* = id.
\end{itemize}\medskip
\ifdefined\hypertarget\hypertarget{sourcepointer-1}{}\fi
\noindent\textit{Source reference.} Original reference 02QL: sources/stacks-chow/pinned-original-source.tex, source lines 1372--1372. Same-chapter pointer into the author environment situation at source line 1372; invoked at proof line None of the counted proof. It is retained only as a navigation pointer; every load-bearing step is inside the delivered span.\par\medskip
\ifdefined\hypertarget\hypertarget{sourcepointer-2}{}\fi
\noindent\textit{Source reference.} Original lemma 02RD: sources/stacks-chow/pinned-original-source.tex, source lines 2287--2287. Same-chapter pointer into the author environment lemma at source line 2287; invoked at proof line 13506 of the counted proof. It is retained only as a navigation pointer; every load-bearing step is inside the delivered span.\par\medskip
\ifdefined\hypertarget\hypertarget{sourcepointer-3}{}\fi
\noindent\textit{Source reference.} Original lemma 067R: sources/stacks-chow/tags.txt, source lines 7980--7980. Cross-chapter standard prerequisite located in the project tag list staged in the local snapshot; invoked at proof line 13516 of the counted proof. It is retained only as a navigation pointer; every load-bearing step is inside the delivered span.\par\medskip
\ifdefined\hypertarget\hypertarget{sourcepointer-4}{}\fi
\noindent\textit{Source reference.} Original lemma 067P: sources/stacks-chow/tags.txt, source lines 7978--7978. Cross-chapter standard prerequisite located in the project tag list staged in the local snapshot; invoked at proof line 13517 of the counted proof. It is retained only as a navigation pointer; every load-bearing step is inside the delivered span.\par\medskip
\ifdefined\hypertarget\hypertarget{sourcepointer-5}{}\fi
\noindent\textit{Source reference.} Original lemma 0FF0: sources/stacks-chow/pinned-original-source.tex, source lines 13412--13412. Same-chapter pointer into the author environment lemma at source line 13412; invoked at proof line 13518 of the counted proof. It is retained only as a navigation pointer; every load-bearing step is inside the delivered span.\par\medskip
\ifdefined\hypertarget\hypertarget{sourcepointer-6}{}\fi
\noindent\textit{Source reference.} Original reference 0F9Z: sources/stacks-chow/pinned-original-source.tex, source lines 6184--6184. Same-chapter pointer into the author environment remark at source line 6184; invoked at proof line 13521 of the counted proof. It is retained only as a navigation pointer; every load-bearing step is inside the delivered span.\par\medskip
\ifdefined\hypertarget\hypertarget{sourcepointer-7}{}\fi
\noindent\textit{Source reference.} Original lemma 0FBK: sources/stacks-chow/pinned-original-source.tex, source lines 11912--11912. Same-chapter pointer into the author environment lemma at source line 11912; invoked at proof line 13522 of the counted proof. It is retained only as a navigation pointer; every load-bearing step is inside the delivered span.\par\medskip
\ifdefined\hypertarget\hypertarget{sourcepointer-8}{}\fi
\noindent\textit{Source reference.} Original lemma 0FF1: sources/stacks-chow/pinned-original-source.tex, source lines 13435--13435. Same-chapter pointer into the author environment lemma at source line 13435; invoked at proof line 13526 of the counted proof. It is retained only as a navigation pointer; every load-bearing step is inside the delivered span.\par\medskip
\begin{lemma}
\label{lemma-lci-gysin-well-defined}
Let $(S, \delta)$ be as in Situation \ref{situation-setup}.
Let $f : X \to Y$ be a local complete intersection morphism
of schemes locally of finite type over $S$.
The bivariant class $f^!$ is independent of the choice of
the factorization $f = g \circ i$ with $g$ smooth (provided
one exists).
\end{lemma}

\begin{proof}
Given a second such factorization $f = g' \circ i'$ we can
consider the smooth morphism $g'' : P \times_Y P' \to Y$, the
immersion $i'' : X \to P \times_Y P'$ and the factorization
$f = g'' \circ i''$. Thus we may assume that we have a diagram
$$
\xymatrix{
& P' \ar[d]^p \ar[rd]^{g'} \\
X \ar[r]^i \ar[ru]^{i'} & P \ar[r]^g & Y
}
$$
where $p$ is a smooth morphism. Then $(g')^* = p^* \circ g^*$
(Lemma \ref{lemma-compose-flat-pullback}) and hence it suffices
to show that $i^! = (i')^! \circ p^*$
in $A^*(X \to P)$. Consider the commutative diagram
$$
\xymatrix{
& X \times_P P' \ar[d]^{\overline{p}} \ar[r]_j & P' \ar[d]^p \\
X \ar[ru]^s \ar[r]^1 & X \ar[r]^i & P
}
$$
where $s =(1, i')$. Then $s$ and $j$ are regular immersions
(by Divisors, Lemma \ref{divisors-lemma-section-smooth-regular-immersion}
and Divisors, Lemma \ref{divisors-lemma-flat-base-change-regular-immersion})
and $i' = j \circ s$. By Lemma \ref{lemma-composition-regular-immersion}
we have $(i')^! = s^! \circ j^!$.
Since the square is cartesian, the bivariant class $j^!$
is the restriction (Remark \ref{remark-restriction-bivariant})
of $i^!$ to $P'$, see Lemma \ref{lemma-construction-gysin}.
Since bivariant classes commute with flat pullbacks
we find $j^! \circ p^* = \overline{p}^* \circ i^!$.
Thus it suffices to show that $s^! \circ \overline{p}^* = \text{id}$
which is done in Lemma \ref{lemma-section-smooth}.
\end{proof}

\end{document}
